\section{Graph Learning for Failure-Impact Prediction}
\label{sec:5}

While the interpretable RM baseline provides explainable quality profiles, complex non-linear failure interactions and non-local cascade propagations benefit significantly from graph representation learning. This section details our Heterogeneous Graph Transformer (HGT) architecture, 16-dimensional edge representation, multi-task prediction heads, dimension-masked loss formulation, ground-truth simulation oracles, and the strict input--label independence guarantee.

\subsection{Ground-Truth Simulation Oracles}
\label{sec:5.1}

To evaluate predictive accuracy prior to deployment without relying on production runtime telemetry, SaG executes discrete-event failure simulations over the raw structural multigraph $G_{\text{structural}}$. We establish a formal three-oracle taxonomy:

\begin{itemize}
    \item \textbf{Cascade Reachability Oracle ($I^*(v)$):} Implemented via \texttt{FaultInjector}, this oracle injects an unexpected node crash at component $v \in V$, propagates cascading failures across dependent topics, brokers, and network links using breadth-first dynamic traversal, and computes the fraction of surviving subscriber feeds severed by the outage. $I^*(v) \in [0, 1]$ serves as the primary continuous target label for training and evaluating GNN predictors.
    \item \textbf{Multi-Metric Composite Oracle ($I_{\text{comp}}(v)$):} Implemented via \texttt{FailureSimulator}, this oracle evaluates a multi-faceted failure impact vector:
    \begin{equation}
    I_{\text{comp}}(v) = 0.35 \cdot \Delta\text{Reachability} + 0.25 \cdot \Delta\text{Fragmentation} + 0.25 \cdot \Delta\text{Throughput} + 0.15 \cdot \Delta\text{FlowDisruption}
    \end{equation}
    $I_{\text{comp}}(v)$ serves as the canonical oracle for architectural quality gate verification.
    \item \textbf{Dynamic Queue-Flow Oracle ($I_{\text{dyn}}(v)$):} Implemented via \texttt{MessageFlowSimulator} using the SimPy discrete-event simulation framework \cite{team2020simpy}, this oracle models message emission rates, stochastic network latencies, broker buffer saturation, and dropped delivery counts under fault injection.
    \item \textbf{Relationship (Edge) Removal Oracle ($I_{\text{edge}}(u,v)$):} Evaluates the systemic impact of severing an individual dependency or communication channel while keeping both endpoint components alive:
    \begin{equation}
    I_{\text{edge}}(u,v) = I_{\text{comp}}(G \setminus \{(u,v)\}) - I_{\text{comp}}(G)
    \end{equation}
\end{itemize}

\subsection{Heterogeneous Graph Transformer Architecture}
\label{sec:5.2}

Because distributed systems comprise heterogeneous entity types (Applications, Libraries, Brokers, Topics, Infrastructure Nodes) and diverse interaction semantics (\texttt{PUBLISHES\_TO}, \texttt{SUBSCRIBES\_TO}, \texttt{RUNS\_ON}, \texttt{CONNECTS\_TO}, \texttt{USES}, etc.), we implement a 3-layer \textbf{Heterogeneous Graph Transformer (HGT)} architecture \cite{hu2020hgt} with hidden dimension $D = 64$ and 4 attention heads.

\subsubsection{Edge Feature Encoding (16-Dimensional)}
\label{sec:5.2.1}

To capture continuous QoS constraints and channel semantics, SaG encodes each directed edge $e = (u, v)$ as a 16-dimensional continuous-categorical vector $e_{uv} \in \mathbb{R}^{16}$:
\begin{itemize}
    \item \textbf{Index 0:} Scalar QoS weight $w(e) \in (0, 1]$, computed via the harmonic mean of reliability, durability, priority, deadline, and blocking multipliers.
    \item \textbf{Index 1:} Normalized path count through edge $e$ in $G_{\text{analysis}}$.
    \item \textbf{Indices 2--8:} 7-bit one-hot encoding for edge relationship types (\texttt{PUBLISHES\_TO}, \texttt{SUBSCRIBES\_TO}, \texttt{ROUTES}, \texttt{RUNS\_ON}, \texttt{CONNECTS\_TO}, \texttt{USES}, \texttt{DEPENDS\_ON}).
    \item \textbf{Indices 9--15:} 7 explicit QoS profile dimensions, non-zero only for \texttt{PUBLISHES\_TO}/\texttt{SUBSCRIBES\_TO} edges (all other edge types receive zeros): (9) Reliability policy score (0.0 = best-effort, 1.0 = reliable); (10) Durability policy score (0.0 = volatile, 0.5 = transient local, 0.6 = transient, 1.0 = persistent); (11) Normalized message priority (0.0/0.33/0.66/1.0 for low/medium/high/urgent); (12) Deadline constraint active flag (0/1); (13) Log deadline $\log_{10}(1 + \text{deadline\_ns} / 10^6)$; (14) Log max blocking time $\log_{10}(1 + \text{max\_blocking\_ms})$; and (15) a QoS heterogeneity flag (1 if the edge's reliability/durability/priority triple deviates from the scenario's modal QoS profile, 0 otherwise).
\end{itemize}

The \texttt{EdgeFeatureEncoder} projects $e_{uv}$ into the hidden space: $e_{uv}' = W_{\text{edge}} e_{uv}$. Prior to relational attention, the edge projection is added directly to the destination node representation: $\tilde{h}_v = h_v + e_{uv}'$.

\subsubsection{Type-Specific Projection and Heterogeneous Message Passing}

For each source node $u$ and target node $v$ with meta-relation $\tau(e) = (\tau(u), \phi(e), \tau(v))$:
\begin{enumerate}
    \item \textbf{Type-Specific Projection:} Node feature vectors $x_v$ (dimensions 19--23 depending on $\tau(v)$) are mapped into $D$-dimensional hidden space:
    \begin{equation}
    h_v^{(0)} = \text{LayerNorm}(\text{GELU}(W_{\tau(v)} x_v))
    \end{equation}
    \item \textbf{Relational Mutual Attention:} Query, Key, and Value projections compute relation-specific attention across heads:
    \begin{equation}
    \text{Attn}(u, e, v) = \underset{\forall u \in \mathcal{N}(v)}{\text{Softmax}}\left( \frac{K(u) W_{\text{att},\phi(e)} Q(\tilde{v})^\top}{\sqrt{D/H}} \right)
    \end{equation}
    \begin{equation}
    \text{Msg}(u, e, v) = V(u) W_{\text{msg},\phi(e)}
    \end{equation}
    \item \textbf{Bidirectional Message Passing:} To capture downstream backpressure and upstream cascading failures, message passing is executed over both forward and reverse relations ($G_{\text{analysis}}$ and $G_{\text{analysis}}^\top$).
    \item \textbf{Residual Aggregation and Layer Normalization:}
    \begin{equation}
    h_v^{(l)} = \text{LayerNorm}\left( h_v^{(l-1)} + \text{Dropout}\left(\sum_{u \in \mathcal{N}(v)} \text{Attn}(u, e, v) \cdot \text{Msg}(u, e, v)\right)\right)
    \end{equation}
\end{enumerate}

\subsection{Multi-Task Prediction Heads and Dimension Masking}
\label{sec:5.3}

From the final node embeddings $h_v^{(L)}$, SaG branches into specialized multi-task prediction heads:
\begin{itemize}
    \item \textbf{Reliability Head:} $\hat{R}(v) = \sigma(\text{MLP}_R(h_v)) \in [0, 1]$
    \item \textbf{Maintainability Head:} $\hat{M}(v) = \sigma(\text{MLP}_M(h_v)) \in [0, 1]$
    \item \textbf{Composite Failure Impact Head:} $\hat{I}^*(v) = \sigma(\text{MLP}_C(h_v \parallel \hat{R}(v) \parallel \hat{M}(v))) \in [0, 1]$
    \item \textbf{Relationship Criticality Head:} $\hat{Q}(u,v) = \sigma(\text{TypedEdgeEncoder}_{\phi(e)}(h_u, h_v, e_{uv})) \in [0, 1]$
\end{itemize}

\subsubsection{Dimension-Masked Loss Formulation}

The joint optimization objective balances regression accuracy, multi-task dimension learning, ranking fidelity, pairwise ordering, and edge prediction:
\begin{equation}
\mathcal{L} = \mathcal{L}_{\text{composite}} + 0.5 \cdot \mathcal{L}_{\text{dimension}} + 0.3 \cdot \mathcal{L}_{\text{rank}} + 0.1 \cdot \mathcal{L}_{\text{pairwise}} + 0.1 \cdot \mathcal{L}_{\text{consistency}} + 0.3 \cdot \mathcal{L}_{\text{edge}}
\end{equation}
where $\mathcal{L}_{\text{composite}} = \text{MSE}(\hat{I}^*(v), I^*(v))$, $\mathcal{L}_{\text{rank}}$ is ListMLE ranking loss \cite{xia2008listwise}, $\mathcal{L}_{\text{pairwise}}$ is margin-ranking loss, and $\mathcal{L}_{\text{consistency}} = \text{MSE}(\hat{I}^*(v), 0.8\hat{R}(v) + 0.2\hat{M}(v))$.

\noindent\textbf{Dimension Masking:} Because dynamic cascade simulation ($I^*(v)$ via \texttt{FaultInjector}) observes runtime failure reachability rather than static code maintainability, maintainability ground truth is unobserved during dynamic simulation. We introduce a boolean dimension mask $m = [m_R, m_M] = [1, 0]$:
\begin{equation}
\mathcal{L}_{\text{dimension}} = \frac{1}{\sum_{d} m_d} \sum_{d \in \{R, M\}} m_d \cdot \text{MSE}(\hat{d}(v), d^*(v))
\end{equation}
This prevents the unmeasured maintainability head from being penalized or driven toward zero by backpropagation.

\subsubsection{Domain-Reweighted Criticality (\texorpdfstring{$Q_{\text{domain}}$}{Qdomain})}

To ground predictions in ISO/IEC 25019 Context of Use ($\vec{\omega} = [q_R, q_M]^\top$), the composite score is evaluated as:
\begin{equation}
Q_{\text{domain}}(v) = q_R \cdot \hat{R}(v) + q_M \cdot M_{\text{static}}(v)
\end{equation}
where $M_{\text{static}}(v)$ is drawn directly from the structural analyzer's maintainability baseline ($y_{\text{rm}}$ maintainability column), combining learned dynamic reliability with static source-code maintainability. The implementation deliberately refuses to emit $Q_{\text{domain}}(v)$ unless the checkpoint's maintainability head actually received real supervision --- under the sole labeler used in every experiment reported here (\texttt{FaultInjector}, $m = [1, 0]$), maintainability is unmeasured, so $Q_{\text{domain}}(v)$ is never populated in these results. Domain reweighting sensitivity is instead evaluated directly against the static RM baseline, reported as a dedicated ablation (Section~\ref{sec:7.3}).

\subsection{Input--Label Independence Guarantee}
\label{sec:5.4}

To guarantee experimental rigor and eliminate data leakage, SaG enforces strict architectural decoupling:
\begin{itemize}
    \item \textbf{Feature Space:} Derived exclusively from $G_{\text{analysis}}$ (using static structural topology, SCA metrics, and declared QoS contracts).
    \item \textbf{Label Space:} Evaluated exclusively on raw $G_{\text{structural}}$ through independent simulation oracles (\texttt{FaultInjector}, \texttt{FailureSimulator}, \texttt{MessageFlowSimulator}).
\end{itemize}
No simulation outputs are ever exposed as input features to the GNN or attribution baseline.
